\documentclass[a4paper,11pt,ja=standard,xelatex]{bxjsarticle}

\usepackage{fontspec}
\usepackage{xeCJK}
\setmonofont{DejaVu Sans Mono}
\setCJKmainfont{Noto Serif CJK JP}
\setCJKsansfont{Noto Sans CJK JP}
\setCJKmonofont{Noto Sans Mono CJK JP}

\geometry{top=24mm,bottom=25mm,left=23mm,right=23mm}
\usepackage{fancyhdr}
\usepackage{lastpage}
\usepackage{longtable}
\usepackage{booktabs}
\usepackage{tabularx}
\usepackage{array}
\usepackage{enumitem}
\usepackage[most]{tcolorbox}
\usepackage[unicode,hidelinks]{hyperref}
\usepackage{xurl}

\setlength{\parindent}{1em}
\setlength{\headheight}{16pt}
\setlength{\parskip}{0.55em}
\setlist[itemize]{topsep=0.2em,itemsep=0.15em,leftmargin=2em}
\setlist[enumerate]{topsep=0.2em,itemsep=0.15em,leftmargin=2em}
\renewcommand{\arraystretch}{1.25}
\setlength{\tabcolsep}{5pt}
\sloppy

\pagestyle{fancy}
\fancyhf{}
\fancyhead[L]{第04章 ソフトウェア構築}
\fancyhead[R]{SWEBOK v4.0a 日本語まとめ}
\fancyfoot[C]{\thepage/\pageref{LastPage}}
\renewcommand{\headrulewidth}{0.4pt}
\renewcommand{\footrulewidth}{0pt}

\newcolumntype{Y}{>{\raggedright\arraybackslash}X}
\newcolumntype{L}[1]{>{\raggedright\arraybackslash}p{#1}}

\newtcolorbox{plainbox}{
  breakable,
  colback=white,
  colframe=black,
  boxrule=0.5pt,
  sharp corners,
  left=5pt,right=5pt,top=5pt,bottom=5pt
}


\newtcolorbox{keybox}[1]{
  breakable,
  colback=white,
  colframe=black,
  boxrule=0.5pt,
  sharp corners,
  left=5pt,right=5pt,top=5pt,bottom=5pt,
  before upper={\textbf{#1}\par\smallskip}
}

\newtcolorbox{todobox}[1]{
  breakable,
  colback=white,
  colframe=black,
  boxrule=0.5pt,
  sharp corners,
  left=5pt,right=5pt,top=5pt,bottom=5pt,
  before upper={\textbf{TODO（図）　#1}\par\smallskip}
}

\newcommand{\eng}[1]{\textsf{#1}}
\newcommand{\topic}[2]{\noindent\textbf{#1}\quad #2\par}

\begin{document}
\thispagestyle{empty}

\begin{center}
  {\LARGE 第04章 ソフトウェア構築}\par
  \vspace{0.8em}
  {\large Software Construction の日本語まとめ}\par
  \vspace{0.6em}
  SWEBOK Guide v4.0a Chapter 04 を初学者向けに再構成\par
  \vspace{1.2em}
  作成日：2026 年 4 月 29 日
\end{center}

\vspace{1em}
\hrule
\vspace{0.8em}

\noindent\textbf{この資料の主張}

ソフトウェア構築は、単にプログラムを書く作業ではない。要求と設計を、実行でき、検証でき、保守できるソフトウェアへ変える中心的な活動である。構築には、コーディング、検証、単体テスト、統合テスト、デバッグ、再利用、統合、依存関係管理、性能調整、開発環境の選択が含まれる。よい構築は、複雑さを抑え、変化を受け入れ、欠陥を早く見つけ、既存資産を適切に使い、標準に沿って進めることで実現する。

\begin{plainbox}
\textbf{この資料の読み方}

専門用語は原則として日本語で示し、必要に応じて英語を括弧内に併記する。英語名は、元資料やツール上の名称を確認したいときの手がかりとして残す。初学者が前提知識なしに読めるように、各節では「何のための概念か」「実務では何を見ればよいか」を重視する。
\end{plainbox}

\vspace{1em}
\tableofcontents

\newpage

\section*{全体概要：この章で学ぶこと}
\addcontentsline{toc}{section}{全体概要：この章で学ぶこと}

この章は、ソフトウェアを実際に作る段階で必要になる知識を扱う。ここでいう「作る」とは、ソースコードを書くことだけではない。詳細設計を補い、テストし、部品を統合し、依存関係を管理し、性能を測り、運用に近い環境からのフィードバックを得ながら、ソフトウェアを完成度の高い形に近づけることである。

\noindent\textbf{到達目標}\quad この章を読み終えると、次のことができるようになる。

\begin{itemize}
  \item ソフトウェア構築が、設計、テスト、品質、構成管理、プロジェクト管理とどう関係するかを説明できる。
  \item 複雑さの最小化、変化への備え、検証しやすさ、再利用、標準適用の意味を説明できる。
  \item コーディング、単体テスト、統合、依存関係管理、継続的統合の実務上の役割を説明できる。
  \item API、例外処理、実行可能モデル、状態機械、ミドルウェア、クラウド、並行処理、性能調整などの構築技術を概観できる。
  \item 開発環境、単体テストツール、性能分析ツール、ローコード/ゼロコード環境の位置づけを説明できる。
\end{itemize}

\noindent\textbf{学習順序}\quad 初学者は、まず「構築とは何か」と「構築の基本原則」を押さえる。その後、「構築をどう管理するか」「現場で何を考えるか」「どの技術を使うか」「どのツールが支えるか」という順で読むとよい。

\begin{todobox}{第04章の全体地図を入れる}
\textbf{画像生成用プロンプト}

白背景、モノクロ、A4 本文幅に収まる横長の学習用図解を作成する。中央に「ソフトウェア構築」を置き、周囲に「基礎」「構築の管理」「実務上の考慮」「構築技術」「構築ツール」を配置する。さらに外側に「設計」「テスト」「品質」「構成管理」「プロジェクト管理」「計算基礎」を配置し、関係を細い矢印で示す。ラベルは日本語、講義ノート風、余白を広めにする。
\end{todobox}

\section*{最初に必要な用語}
\addcontentsline{toc}{section}{最初に必要な用語}

\begin{longtable}{L{0.23\linewidth}L{0.69\linewidth}}
\toprule
\textbf{用語} & \textbf{初学者向けの説明} \\
\midrule
ソフトウェア構築 & ソフトウェアを詳細に作成・保守する活動である。コーディング、検証、単体テスト、統合テスト、デバッグなどを含む。 \\
コーディング & プログラミング言語などを使い、設計や要求を実行可能な形に書き表すこと。 \\
検証 & 作ったものが、仕様、設計、標準、期待された振る舞いを満たしているか確認すること。 \\
単体テスト & クラス、関数、モジュールなど、個別の部品を単独で確認するテスト。 \\
統合テスト & 複数の部品を組み合わせたとき、相互作用が正しく働くかを確認するテスト。 \\
デバッグ & 失敗の原因となる欠陥を見つけ、原因を理解し、修正する作業。 \\
依存関係 & 自分のソフトウェアが利用する外部ライブラリ、フレームワーク、部品、サービスなどとの関係。 \\
継続的統合 & 変更を頻繁に統合し、自動ビルドと自動テストで早く問題を発見する実践。 \\
再利用 & 既存の部品、コード、ライブラリ、サービス、設計パターンなどを別の問題解決に使うこと。 \\
API & アプリケーションプログラミングインタフェースのことである。ライブラリやサービスを利用するために公開された操作、引数、意味の取り決めである。 \\
\bottomrule
\end{longtable}

\section*{略語}
\addcontentsline{toc}{section}{略語}

\begin{longtable}{L{0.16\linewidth}L{0.32\linewidth}L{0.44\linewidth}}
\toprule
\textbf{略語} & \textbf{日本語} & \textbf{説明} \\
\midrule
API & アプリケーションプログラミングインタフェース & ライブラリ、サービス、フレームワークを使うための公開された入り口。 \\
CI & 継続的統合 & 変更を頻繁に統合し、自動ビルドと自動テストで確認する実践。 \\
COTS & 商用既製品 & 自作ではなく、市販または既存の製品として利用するソフトウェア部品。 \\
DSL & ドメイン固有言語 & 特定領域の問題を表しやすいように作られた言語。 \\
IDE & 統合開発環境 & 編集、ビルド、テスト、デバッグなどを統合した開発環境。 \\
MDA & モデル駆動アーキテクチャ & モデルから実装へ変換する考え方を中心にした開発方式。 \\
PIM & プラットフォーム非依存モデル & 特定の実装環境に依存しない解決策のモデル。 \\
PSM & プラットフォーム固有モデル & 実装環境の詳細を含むモデル。 \\
TDD & テスト駆動開発 & テストを先に書き、そのテストを通すようにコードを書く開発スタイル。 \\
WYSIWYG & 見たままが得られる方式 & 画面上の見た目が成果物の見た目に近い編集方式。 \\
\bottomrule
\end{longtable}

\section*{導入：ソフトウェア構築とは何か}
\addcontentsline{toc}{section}{導入：ソフトウェア構築とは何か}

ソフトウェア構築とは、ソフトウェアを詳細に作成し、保守する活動である。具体的には、コーディング、検証、単体テスト、統合テスト、デバッグを含む。構築は、設計書の内容をただ機械的に写す作業ではない。多くの現場では、詳細な設計判断の一部が構築中に行われる。たとえば、どのデータ構造を使うか、エラーをどこで検出するか、処理をどの関数に分けるか、どのライブラリを使うかは、構築段階で決まることが多い。

この章の知識領域は、他の知識領域と強く結びつく。設計の成果物を入力として使い、テストに渡す成果物を作るため、ソフトウェア設計とソフトウェアテストとの関係が特に強い。また、構築ではソースファイル、テストケース、設定ファイル、文書など多数の構成品目が生まれるため、ソフトウェア構成管理とも密接に関係する。品質、プロジェクト管理、計算基礎とも関係が深い。

\begin{keybox}{重要}
ソフトウェア構築は、実装だけでなく「実装を検証可能にし、変更可能にし、統合可能にし、保守可能にする活動」である。したがって、構築の品質は、完成したソフトウェアの品質だけでなく、後続の変更や運用のしやすさにも影響する。
\end{keybox}

\begin{todobox}{構築と他知識領域の関係図を入れる}
\textbf{画像生成用プロンプト}

白背景、モノクロ、A4 本文幅の関係図を作成する。中央に「ソフトウェア構築」。左に「ソフトウェア設計」、右に「ソフトウェアテスト」、下に「構成管理」「品質」「プロジェクト管理」「計算基礎」を置く。設計から構築へ「設計成果物」、構築からテストへ「コード・単体テスト・統合結果」と矢印を付ける。講義ノート風、細線、日本語ラベル。
\end{todobox}

\section{ソフトウェア構築の基礎}

ソフトウェア構築の基礎は、五つの考え方で整理できる。第一に複雑さを最小化すること、第二に変化を予期し受け入れること、第三に検証しやすく作ること、第四に資産を再利用すること、第五に標準を適用することである。これらは設計にも当てはまるが、構築では特にコード、テスト、部品、ツールに直接表れる。

\begin{todobox}{構築の五つの基礎概念を示す図を入れる}
\textbf{画像生成用プロンプト}

白背景、モノクロ、A4 本文幅の放射状図解を作成する。中央に「よいソフトウェア構築」を置き、五方向に「複雑さの最小化」「変化を予期し受け入れる」「検証しやすく構築する」「資産を再利用する」「標準を適用する」を配置する。各項目の下に短い説明を一行ずつ入れる。講義ノート風、細線、余白を広めにする。
\end{todobox}

\subsection{複雑さの最小化}

人間の作業記憶には限界がある。複雑な構造や大量の詳細を長時間正確に保持することは難しい。この限界は、ソフトウェア構築の中心課題である。コードが複雑になると、理解しにくく、テストしにくく、変更しにくくなる。したがって、構築では「賢そうなコード」よりも「単純で読みやすいコード」を優先する必要がある。

複雑さには複数の種類がある。条件分岐やループが多いコードは、実行経路が多くなる。依存関係が多いコードは、変更の影響範囲が広くなる。名前がわかりにくいコードは、読者の理解を妨げる。ビルド手順が複雑であれば、構築そのものが不安定になる。

循環的複雑度は、コードの制御フローの複雑さを測る代表的な指標である。大まかには、独立した実行経路の数を表す。値が大きいほど、理解やテストが難しくなる傾向がある。ただし、指標は判断の補助であり、指標だけでよいコードかどうかを決めてはならない。

複雑さを下げる実務上の手段には、命名規則、適切な関数分割、モジュール化、重複の排除、標準的な制御構造の利用、読みやすいレイアウト、レビュー、静的解析、単体テストがある。

\begin{keybox}{見分け方}
複雑なコードは、説明に時間がかかる。単純なコードは、名前、構造、流れから意図が読み取れる。構築では「動くか」だけでなく「次の人が読めるか」を基準にする必要がある。
\end{keybox}

\subsection{変化を予期し受け入れる}

多くのソフトウェアは時間とともに変わる。要求が変わることもあれば、利用者数、実行環境、ライブラリ、法規制、ビジネス方針が変わることもある。構築で変化を予期するとは、将来の変更が入ったときに、局所的な修正で済むように作ることである。

変化を予期するためには、拡張しやすい構造、明確なインタフェース、設定による切り替え、再利用可能な部品、テストによる安全網が必要である。一方で、まだ必要でない過剰な柔軟性を作り込むと、逆に複雑さが増える。したがって、予期すべき変化を見極め、過剰設計を避けることが重要である。

現代の開発では、変化を避けるのではなく受け入れる実践が広く使われる。アジャイル開発、DevOps、継続的デリバリ、継続的デプロイは、短い周期で変更を入れ、テストし、提供し、運用から学ぶための仕組みである。

\subsection{検証しやすく構築する}

検証しやすく構築するとは、コードを書いた人、テスト担当者、利用者が欠陥を見つけやすいようにソフトウェアを作ることである。これは、後でテストを追加すればよいという意味ではない。構築中から、テスト可能性を意識した設計と実装が必要である。

具体的には、単体テストを書きやすい関数分割、依存関係を差し替えやすい構造、複雑な言語機能の制限、ログによる挙動記録、コーディング標準、コードレビュー、入力値の検査、異常系の明示が役立つ。

検証しやすいコードでは、失敗が早く発見される。失敗が早く見つかるほど、原因に近い場所で修正できるため、修正費用を下げやすい。

\subsection{資産の再利用}

再利用とは、既存の資産を別の問題解決に使うことである。構築で再利用される資産には、フレームワーク、ライブラリ、モジュール、コンポーネント、ソースコード、商用既製品、クラウドサービスなどがある。

再利用には二つの側面がある。第一は、再利用されることを意識して作る「再利用のための構築」である。第二は、既存の資産を使って新しい解決策を作る「再利用による構築」である。前者では、汎用性、可変性、文書化、テストが重要になる。後者では、選定、評価、統合、ライセンス、保守、脆弱性管理が重要になる。

\begin{keybox}{注意}
再利用は必ず得とは限らない。再利用した部品に脆弱性、ライセンス問題、保守停止、過剰な依存があれば、後から大きな費用が発生する。再利用するものは、品質と責任範囲を確認する必要がある。
\end{keybox}

\subsection{構築における標準の適用}

標準は、構築を効率化し、品質を安定させ、複数人の作業をそろえるために使う。標準には外部標準と内部標準がある。外部標準は、業界団体、国際標準化機関、プラットフォーム提供者などが定める。内部標準は、組織やプロジェクトが定める。

構築に直接関係する標準には、文書形式、プログラミング言語、コーディング規約、例外処理方針、プラットフォームのインタフェース、ツール記法、UML などの表記法がある。特にセキュリティの観点では、許可する言語機能やライブラリ、禁止する危険な書き方を明確にすることが重要である。

\section{構築の管理}

構築は個人の作業に見えやすいが、実際には管理対象が多い。どのライフサイクルで進めるか、どの順で部品を作るか、いつ統合するか、何を測るか、どの依存関係を許すかによって、品質と進捗は大きく変わる。

\subsection{ライフサイクルモデルにおける構築}

ライフサイクルモデルによって、構築の位置づけは異なる。ウォーターフォール型や段階的提供型では、要求や設計がある程度完了した後に構築が始まる。この場合、構築は主にコーディングとデバッグとして見られやすい。

一方、進化的プロトタイピングやアジャイル開発では、要求、設計、構築、テストが重なり合う。短い反復の中で、必要な要求を具体化し、設計し、コードを書き、テストする。継続的デリバリや継続的デプロイでは、コード、テスト、提供、配置がさらに密接に結びつき、自動化されたパイプラインで実行される。

\begin{todobox}{ライフサイクル別の構築位置を比較する図を入れる}
\textbf{画像生成用プロンプト}

白背景、モノクロ、A4 本文幅の比較図を作成する。上段に「ウォーターフォール型」として「要求」「設計」「構築」「テスト」「運用」を直線で並べる。中段に「アジャイル型」として小さな反復サイクルを複数描き、各サイクルに「要求詳細化」「設計」「構築」「テスト」を入れる。下段に「継続的デリバリ/デプロイ」として「コード」「自動ビルド」「自動テスト」「配置」「フィードバック」を循環で示す。日本語ラベル、講義ノート風。
\end{todobox}

\subsection{構築計画}

構築計画では、構築方法、部品の作成順序、統合順序、統合戦略、品質管理プロセス、担当割り当て、必要なツールや環境を決める。構築方法の選択は、複雑さの低減、変化への対応、検証しやすさに影響する。

たとえば、最初から全体を一度に統合する方針を選べば、初期の計画は単純に見えるが、問題発見が遅れる可能性がある。小さな部品を段階的に統合する方針を選べば、統合のための足場やテスト環境が必要になるが、問題を早く見つけやすい。

\subsection{構築測定}

構築では、多くの活動や成果物を測定できる。開発したコード量、変更したコード量、再利用したコード量、削除したコード量、コード複雑度、コードレビューで見つかった欠陥数、欠陥発見率、欠陥修正率、作業工数、予定との差分などである。

測定の目的は、管理と改善である。数値を集めるだけでは意味がない。たとえば、複雑度が高い箇所を重点的にレビューする、欠陥修正率から計画を見直す、再利用量から開発効率を評価する、といった使い方が必要である。

\subsection{依存関係の管理}

現代のソフトウェアは、多くの外部ライブラリやフレームワークに依存する。Java の Maven、JavaScript の NPM などのパッケージ管理ツールは、依存関係の導入、更新、削除を自動化する。しかし、依存関係は便利である一方、リスクも運ぶ。

依存関係の集合は、供給網のようなネットワークを作る。直接使っているライブラリが、さらに別のライブラリに依存していることも多い。どこかに脆弱性、ライセンス衝突、保守停止、悪意あるコードがあれば、自分の製品にも影響する。

依存関係管理では、不要な依存を避ける、信頼できる依存だけを使う、ライセンスを確認する、脆弱性情報を監視する、バージョン更新を計画する、ロックファイルやソフトウェア部品表を活用する、といった実践が重要である。

\begin{todobox}{依存関係の供給網を示す図を入れる}
\textbf{画像生成用プロンプト}

白背景、モノクロ、A4 本文幅のネットワーク図を作成する。左に「自分の製品」を置き、そこから「直接依存 A」「直接依存 B」「直接依存 C」へ矢印を伸ばす。各直接依存からさらに「間接依存」へ枝分かれさせる。リスクとして「脆弱性」「ライセンス衝突」「保守停止」「不審な依存」を吹き出しで示す。日本語ラベル、講義ノート風。
\end{todobox}

\section{実務上の考慮}

構築は、現実の制約の中で進む。要求は変わり、設計は不完全で、利用できる人員や時間には限りがあり、外部サービスやライブラリにも制約がある。そのため、構築は理論だけでなく、実務的な判断が重要な活動である。

\subsection{構築設計}

構築段階でも設計は行われる。要求分析や上位設計で決まっていない詳細を、構築中に決める必要があるためである。たとえば、アルゴリズム、データ構造、関数やクラスの責務、インタフェース、例外の扱い、ログの粒度などは、構築中に具体化されることが多い。

建築現場で、図面に書かれていない細部を現場で調整する必要があるのと同じで、ソフトウェア構築でも詳細の調整が必要である。ただし、場当たり的に決めるのではなく、設計原則、標準、テスト容易性、変更容易性に基づいて判断する必要がある。

\subsection{構築言語}

構築言語とは、人間がコンピュータに実行可能な解決策を指定するための言語である。狭い意味のプログラミング言語だけでなく、設定言語、ツールキット言語、スクリプト言語、ドメイン固有言語も含む。

\begin{longtable}{L{0.24\linewidth}L{0.66\linewidth}}
\toprule
\textbf{構築言語の種類} & \textbf{説明} \\
\midrule
設定言語 & 限られた選択肢から設定を選び、ソフトウェアの振る舞いや導入形態を変える言語である。設定ファイルなどが該当する。 \\
ツールキット言語 & 特定の部品群やツールキットを組み合わせてアプリケーションを作るための言語である。 \\
スクリプト言語 & 自動化、処理の連結、簡易アプリケーションの構築によく使われる言語である。 \\
プログラミング言語 & 汎用的で柔軟な構築言語である。C/C++、Java、Python、C\# などが例である。 \\
ドメイン固有言語 & 特定の業務領域や技術領域の概念を直接表す言語である。汎用言語よりも領域内の問題を簡潔に表せることがある。 \\
\bottomrule
\end{longtable}

構築言語の表記法には、テキストを中心にする言語的表記、数学的に厳密な形式的表記、図形やアイコンを中心にする視覚的表記がある。言語選択は、性能、信頼性、移植性、セキュリティ、開発者の学習コストに影響する。

\begin{todobox}{構築言語の分類図を入れる}
\textbf{画像生成用プロンプト}

白背景、モノクロ、A4 本文幅の階層図を作成する。最上位に「構築言語」。下に「設定言語」「ツールキット言語」「スクリプト言語」「プログラミング言語」「ドメイン固有言語」を並べる。さらに右側に「言語的表記」「形式的表記」「視覚的表記」という別軸を置き、例として C/C++、Java、Event-B、MATLAB、GUI ビルダーを配置する。日本語、講義ノート風。
\end{todobox}

\subsection{コーディング}

コーディングでは、ソースコードを理解しやすく、変更しやすく、検証しやすく書くことが重要である。単に動作するコードを書くだけでは不十分である。コードは、将来の読者に意図を伝える文書でもある。

構築におけるコーディングの主な考慮事項は次のとおりである。

\begin{itemize}
  \item 命名規則、字下げ、レイアウトにより、理解しやすいソースコードを作る。
  \item クラス、列挙型、変数、名前付き定数などを適切に使う。
  \item 条件分岐、ループ、例外処理などの制御構造を過度に複雑にしない。
  \item 予期できるエラーと例外的なエラーを区別して扱う。
  \item バッファオーバーフロー、配列範囲外参照、入力検証不足など、コードレベルのセキュリティ問題を防ぐ。
  \item スレッド、データベースロック、ファイル、ネットワーク接続などの資源を適切に扱う。
  \item コードを関数、クラス、パッケージなどに整理する。
  \item 必要な文書化とコメントを行う。
  \item 測定に基づいて性能を調整する。
\end{itemize}

\begin{keybox}{コメントの考え方}
コメントは、コードの日本語訳ではない。よいコメントは、コードだけでは読み取りにくい理由、前提、制約、例外、意図を説明する。名前と構造で伝えられることは、名前と構造で伝える方がよい。
\end{keybox}

\subsection{構築テスト}

構築では、主に単体テストと統合テストが行われる。これらは、多くの場合、コードを書いたソフトウェア技術者自身が実行する。構築テストの目的は、欠陥が混入してから発見されるまでの時間差を短くし、修正費用を下げることである。

単体テストは、関数、クラス、モジュールなどを単独で確認する。統合テストは、それらを組み合わせたときに正しく相互作用するかを確認する。構築テストは、通常、システムテスト、アルファテスト、ベータテスト、ストレステスト、使用性テストなどの広い範囲のテストまでは含まない。

テストケースは、コードを書いた後に作られることもあるが、コードを書く前に作られることもある。テストを先に書く方法は、テストファーストプログラミングまたはテスト駆動開発として扱われる。

\subsection{構築における再利用}

構築における再利用には、再利用のための構築と、再利用による構築がある。再利用のための構築では、将来別の場所で使えるように、部品を作る。可変性を設定で変えられるようにする、使い方を文書化する、単体テストを整備する、公開方法を整えるといった作業が必要である。

再利用による構築では、既存の部品を使って新しいソフトウェアを作る。ライブラリ、フレームワーク、オープンソース部品、商用既製品、クラウドサービスが代表例である。クラウドサービスを利用する場合、認証、メッセージング、ストレージなどを外部のサービスへ委任できる。このような形は、バックエンドサービスとしての利用とも呼ばれる。

再利用では、部品を選び、評価し、統合し、テストし、利用情報を記録する必要がある。ライセンス、保守状況、脆弱性、利用条件を確認しない再利用は、後のリスクになる。

\subsection{構築品質}

構築品質は、コードやその近くにある成果物の品質に焦点を当てる。要求や上位設計に関する品質活動も重要だが、構築品質では、詳細設計、コード、単体テスト、統合テスト、ビルド、静的解析、レビューが中心になる。

構築品質を高める主な技法には、単体テスト、統合テスト、テストファースト開発、表明、防御的プログラミング、デバッグ、コードインスペクション、技術レビュー、セキュリティレビュー、静的解析がある。

特にセキュリティでは、よく知られた脆弱性パターンを避ける必要がある。入力検証不足、境界チェック不足、例外の握りつぶし、危険なライブラリの利用、認可確認漏れなどは、構築段階で入りやすい問題である。

\subsection{統合}

統合とは、個別に作られた関数、クラス、部品、サブシステムを一つのシステムへ組み合わせることである。さらに、他のソフトウェアやハードウェアと接続することも統合に含まれる。

統合には、段階的統合と一括統合がある。一括統合は、すべての部品ができてからまとめて結合する方法であり、ビッグバン統合とも呼ばれる。問題発見が遅れやすい。段階的統合は、小さな単位から順に組み合わせる方法であり、エラーの場所を特定しやすく、進捗を把握しやすい。

段階的統合では、まだ存在しない部品の代役として、スタブ、ドライバ、モックオブジェクトなどの足場が必要になることがある。現代では、継続的統合によって、変更のたびに自動的にビルドとテストを行うことが広く使われる。

\begin{todobox}{段階的統合と継続的統合の流れを入れる}
\textbf{画像生成用プロンプト}

白背景、モノクロ、A4 本文幅の工程図を作成する。左側に「部品 A」「部品 B」「部品 C」を個別に描き、中央で「段階的統合」によって順に結合される矢印を示す。右側に「継続的統合パイプライン」として「コミット」「自動ビルド」「単体テスト」「統合テスト」「結果通知」を縦または横に並べる。スタブ、ドライバ、モックを足場として小さく表示する。日本語、講義ノート風。
\end{todobox}

\subsection{クロスプラットフォーム開発と移行}

モバイルアプリケーションなどは、特定のプラットフォームに強く依存することが多い。たとえば、iOS と Android では、開発環境、API、利用できる機能、利用者体験が異なる。各プラットフォーム向けに別々に開発すると、費用と時間が増え、実装差も生まれやすい。

クロスプラットフォーム開発は、共通の言語やフレームワークで作り、複数のプラットフォームに出力する方法である。方式としては、共通言語から各プラットフォーム向けのネイティブアプリを生成する方法と、Web 技術で作った部分をネイティブのコンテナで包むハイブリッド方式がある。

既存アプリケーションを別のプラットフォームへ移す場合は、移行が必要になる。移行では、プログラミング言語の違い、API の違い、データ形式、テスト環境、利用者体験の差を考慮する必要がある。

\section{構築技術}

構築技術は、実際にコードや実行可能なソフトウェアを作るための具体的な方法である。この章では、API、オブジェクト指向、総称、表明、例外処理、実行可能モデル、状態機械、設定、国際化、構文解析、並行処理、ミドルウェア、分散・クラウド、異種システム、性能、標準、テストファースト、フィードバックループを扱う。

\subsection{API の設計と利用}

API は、ライブラリやフレームワーク、サービスを利用するために公開された操作の集合である。API は、関数名や引数の型といった署名だけでなく、その操作が何を意味し、どのような効果を持つかという意味も含む。

よい API は、学びやすく、覚えやすく、読みやすいコードにつながり、誤用しにくく、拡張しやすく、必要な機能を十分に備え、後方互換性を保ちやすい。広く使われるライブラリやサービスでは、実装よりも API の方が長く生きることが多いため、API は安定性が重要である。

Web API では、RESTful API や OpenAPI のような標準的な記述がよく使われる。API ファーストとは、アプリケーションの内部実装より先に API の契約を設計し、その契約に基づいてサーバ側とクライアント側を作る考え方である。

\begin{todobox}{API ファーストの流れを入れる}
\textbf{画像生成用プロンプト}

白背景、モノクロ、A4 本文幅のフロー図を作成する。「API 契約を設計」「OpenAPI などで記述」「サーバ側コード生成」「クライアント側コード生成」「実装」「テスト」の箱を左から右へ並べる。API 契約の下に「操作」「引数」「戻り値」「意味」「エラー」を小さく示す。日本語、講義ノート風。
\end{todobox}

\subsection{オブジェクト指向の実行時問題}

オブジェクト指向言語には、実行時に振る舞いを決める仕組みがある。代表例は多態性と反射である。

多態性とは、一般的な操作を、実際の具体的なオブジェクトの種類に応じて実行時に決められる性質である。たとえば、同じ「描画する」という操作でも、円、四角形、文字列で実行される処理が異なる。このように実行時に呼び出し先が決まることを動的束縛という。

反射とは、プログラムが実行時に自分自身の構造や振る舞いを調べたり、変更したりできる性質である。たとえば、クラス名やメソッド名を文字列として扱い、実行時にオブジェクトを生成したり、メソッドを呼び出したりできる。

これらは柔軟性を高めるが、理解やテストを難しくすることがある。利用する場合は、意図を明確にし、過度に動的な構造を避ける必要がある。

\subsection{パラメータ化、テンプレート、総称}

パラメータ化された型は、型やクラスを定義するときに、利用する型の一部を後から与えられるようにする仕組みである。Java では総称、C++ ではテンプレートと呼ばれる。

たとえば、整数のリスト、文字列のリスト、顧客オブジェクトのリストを別々に作るのではなく、「任意の型 T のリスト」として定義し、利用時に T を指定できる。これにより、型安全性を保ちながら再利用しやすい部品を作れる。

\subsection{表明、契約による設計、防御的プログラミング}

表明とは、プログラム中に置かれる実行可能な条件式である。プログラムが想定する前提や不変条件が成り立っているかを実行時に確認する。開発時には欠陥を早く見つける助けになる。

契約による設計とは、各関数やクラスについて、呼び出し前に満たすべき事前条件、呼び出し後に保証する事後条件、常に守られる不変条件を明確にする考え方である。これにより、部品同士の責任範囲が明確になる。

防御的プログラミングとは、不正な入力や想定外の状態から関数を守るために、入力値を検査し、異常な場合の処理を決めることである。防御的に書くことで、欠陥の影響を局所化しやすい。

\begin{todobox}{表明・契約・防御的プログラミングの関係図を入れる}
\textbf{画像生成用プロンプト}

白背景、モノクロ、A4 本文幅の概念図を作成する。中央に「関数またはクラス」を置く。左から「事前条件」、右へ「事後条件」、内部に「不変条件」を示す。周囲に「表明」と「防御的入力検査」を配置し、どの条件を実行時に確認するかを矢印で示す。日本語、講義ノート風。
\end{todobox}

\subsection{エラー処理、例外処理、フォールトトレランス}

エラー処理の方法は、正しさ、頑健性、信頼性に影響する。エラー処理には、警告ログを出す、エラーコードを返す、中立値を返す、次の有効データを使う、処理を停止するなど、複数の方法がある。どれを使うかは、利用者への影響、データの重要性、安全性、復旧可能性によって変わる。

例外処理は、エラーや例外的な出来事を検出し、処理する仕組みである。基本的には、例外を送出し、対応する処理ブロックで捕捉する。例外メッセージには、原因理解に必要な情報を含めるべきである。空の捕捉ブロックで例外を握りつぶすと、問題の発見が遅れる。

フォールトトレランスとは、エラーを検出し、回復するか、影響を閉じ込めることで、ソフトウェアの信頼性を高める技法群である。再試行、補助コード、投票アルゴリズム、値の代替などがある。

\subsection{実行可能モデル}

実行可能モデルとは、特定のプログラミング言語や実装構造の詳細から離れて、実行できる形で振る舞いを表すモデルである。実行可能 UML などが例である。

モデル駆動アーキテクチャでは、プラットフォーム非依存モデルとプラットフォーム固有モデルを区別する。プラットフォーム非依存モデルは、実装環境に依存しない解決策を表す。プラットフォーム固有モデルは、特定のハードウェアやソフトウェア環境の詳細を含む。変換器は、非依存モデルとプラットフォーム情報を組み合わせ、実装を生成する。

\begin{todobox}{実行可能モデル、PIM、PSM の関係図を入れる}
\textbf{画像生成用プロンプト}

白背景、モノクロ、A4 本文幅の変換図を作成する。左に「プラットフォーム非依存モデル PIM」、中央に「変換器」、右に「プラットフォーム固有モデル PSM」、さらに右に「実装」を並べる。下に「対象プラットフォーム情報」を置き、変換器へ矢印を引く。実行可能モデルは PIM の一例として示す。日本語、講義ノート風。
\end{todobox}

\subsection{状態に基づく構築技法と表駆動構築技法}

状態に基づくプログラミングは、有限状態機械を使ってプログラムの振る舞いを表す技法である。状態、イベント、遷移、動作を明確にすることで、複雑な分岐を整理できる。状態遷移図は、仕様、実装、デバッグ、文書化で使える。

表駆動方式は、複雑な \texttt{if} 文や \texttt{case} 文の代わりに、表を使って情報を表す方法である。条件と処理の対応を表にまとめると、ロジックが見通しやすく、変更もしやすくなることがある。表駆動では、何を表に入れるか、どのように効率よく表を参照するかが重要である。

\begin{todobox}{状態機械と表駆動の比較図を入れる}
\textbf{画像生成用プロンプト}

白背景、モノクロ、A4 本文幅の左右比較図を作成する。左に「状態に基づく構築」として、状態 A、状態 B、状態 C とイベントによる遷移矢印を描く。右に「表駆動構築」として、条件列、動作列、次状態列を持つ表を描く。中央に「複雑な分岐を整理する」という共通目的を置く。日本語、講義ノート風。
\end{todobox}

\subsection{実行時設定と国際化}

実行時設定とは、プログラムの値や動作を、実行時に設定ファイルや環境変数などから読み取って決めることである。これにより、コードを変更せずに動作を変えられる。設定を使うことで柔軟性は高まるが、設定値の検証、既定値、変更履歴、機密情報の扱いが重要になる。

国際化とは、複数の言語や地域に対応できるようにソフトウェアを準備することである。これに対応して、特定の地域や言語に合わせる作業を地域化という。画面の文言、エラーメッセージ、日付形式、通貨、文字コード、並び順、右から左へ書く言語などを考慮する必要がある。

国際化を後から追加すると、コードのあちこちに埋め込まれた文字列や地域依存の処理を修正する必要がある。構築時から、文字列を外部化し、文字コードとロケールを意識することが重要である。

\subsection{文法に基づく入力処理}

文法に基づく入力処理は、入力を字句や構文の規則に従って解析する技法である。代表的な例は、プログラミング言語、設定ファイル、数式、検索式、通信プロトコルの解析である。

構文解析では、入力のトークン列から解析木または構文木と呼ばれるデータ構造を作る。解析木は、その入力がどのような構造を持つかを表す。構文解析の後、意味検査やシンボル表の確認を行い、計算処理や変換処理につなげる。

文法に基づく入力処理は、入力検証にも重要である。曖昧な文字列処理で入力を解釈すると、誤解析やセキュリティ問題につながる可能性がある。

\subsection{並行処理の基本部品}

並行処理では、複数の処理が同時または並行して進む。共有資源に同時にアクセスすると競合が起こるため、同期の仕組みが必要である。代表的な基本部品には、セマフォ、モニタ、ミューテックスがある。

セマフォは、複数のプロセスやスレッドが共有資源へアクセスする数を制御するための抽象である。モニタは、共有変数とその操作をひとまとまりにし、同時に一つの処理だけが内部で動くようにする抽象データ型である。ミューテックスは、共有資源への排他的アクセスを実現するための仕組みである。

並行処理の欠陥は再現しにくいことが多い。デッドロック、競合状態、ライブロック、可視性の問題などを避けるには、設計段階から同期方針を明確にし、テストとレビューを組み合わせる必要がある。

\subsection{ミドルウェア}

ミドルウェアとは、オペレーティングシステムより上、アプリケーションより下で、アプリケーションに共通サービスを提供するソフトウェアである。通信、メッセージング、永続化、トランザクション、名前解決、位置透過性などを提供することがある。

ミドルウェアは、部品同士を接続する役割を持つ。メッセージ指向ミドルウェアは、複数のアプリケーション間のメッセージ交換を支援する。企業サービスバスは、サービス指向の通信や連携を支える基盤として使われる。

\subsection{分散およびクラウド基盤ソフトウェアの構築方法}

分散システムは、物理的に離れた複数の計算機がネットワークで接続され、利用者に一体の資源や機能を提供するシステムである。分散ソフトウェアの構築では、並列性、通信、遅延、部分故障、整合性、セキュリティが問題になる。

分散プログラミングには、クライアントサーバ、三層アーキテクチャ、多層アーキテクチャ、分散オブジェクト、疎結合、密結合などの形がある。

クラウド基盤ソフトウェアでは、マイクロサービス、コンテナ、API ゲートウェイ、サービス登録と発見、監視、自動スケーリングなどを考慮する。従来の分散トランザクションでは ACID 特性を重視することが多い。一方、マイクロサービスでは、すべてを強い整合性で統一するのが難しく、SAGA などによる結果整合性を使うことがある。

\begin{todobox}{分散・クラウド構築の概念図を入れる}
\textbf{画像生成用プロンプト}

白背景、モノクロ、A4 本文幅の構成図を作成する。利用者から「API ゲートウェイ」へ矢印を描き、その先に複数の「マイクロサービス」を配置する。各サービスを「コンテナ」で囲み、下に「クラウド基盤」を置く。横に「サービス登録・発見」「監視」「メッセージング」「データストア」を配置する。注記として「部分故障」「通信遅延」「結果整合性」を入れる。日本語、講義ノート風。
\end{todobox}

\subsection{異種システムの構築}

異種システムとは、種類の異なる計算単位が組み合わさったシステムである。GPU、デジタル信号処理プロセッサ、マイクロコントローラ、周辺プロセッサなどが含まれる。組込みシステムは、異種システムであることが多い。

異種システムでは、ハードウェアとソフトウェアを同時に設計する必要がある。これをハードウェア/ソフトウェア協調設計という。複数の仕様言語、複数のシミュレーション環境、インタフェースの整合、共同検証が課題になる。FPGA や ASIC を使ったハードウェア側の試作・実装と、ソフトウェア側の低水準言語への変換が並行して進むことがある。

\subsection{性能分析とチューニング}

性能は、アーキテクチャ、詳細設計、データ構造、アルゴリズム、コードの書き方に影響される。性能分析は、プログラムの実行時情報を集め、時間がかかっている箇所、頻繁に実行される箇所、資源を多く使う箇所を見つける活動である。

コードチューニングは、コードレベルで性能を改善することである。論理式、ループ、データ変換、関数呼び出し、メモリアクセスなどを見直す。重要なのは、推測ではなく測定に基づいて行うことである。測定せずに最適化すると、本当に遅い場所ではない箇所を複雑にしてしまう危険がある。

\begin{keybox}{性能改善の原則}
まず測定し、ボトルネックを特定し、小さく変更し、再測定する。性能のために可読性を大きく下げる場合は、その理由と測定結果を残す必要がある。
\end{keybox}

\begin{todobox}{性能分析とチューニングの流れを入れる}
\textbf{画像生成用プロンプト}

白背景、モノクロ、A4 本文幅の循環図を作成する。「性能要求の確認」「計測」「ボトルネック特定」「改善案」「実装変更」「再計測」を円環で示す。中央に「推測ではなく測定」と書く。横に「実行時間」「メモリ」「I/O」「ネットワーク」の測定対象を小さく配置する。日本語、講義ノート風。
\end{todobox}

\subsection{プラットフォーム標準}

プラットフォーム標準は、互換性のある環境でアプリケーションを変更なしに実行しやすくするための標準である。標準サービスや API を定めることで、開発者は特定環境に過度に依存しないアプリケーションを作りやすくなる。

代表例には、エンタープライズ向けの Jakarta Enterprise Edition、Unix 系オペレーティングシステムのインタフェース標準である POSIX、Web アプリケーションの標準を支える HTML5 がある。

標準に従うことは、移植性、相互運用性、保守性を高める。一方で、標準だけでは解決できないプラットフォーム固有の差もあるため、現実の環境での検証が必要である。

\subsection{テストファーストプログラミング}

テストファーストプログラミングは、コードを書く前にテストケースを書く開発スタイルである。現在のコードベースでそのテストは失敗する。その後、テストが通るようにコードを書く。テストが通ったら、必要に応じてリファクタリングし、設計を整える。

この方法は、テスト駆動開発としても知られる。テストを先に書くことで、プログラマは要求と設計を先に考える。これにより、要求や設計の問題が早く表面化する。欠陥も早く発見され、修正しやすい。

\begin{todobox}{テスト駆動開発のサイクル図を入れる}
\textbf{画像生成用プロンプト}

白背景、モノクロ、A4 本文幅の循環図を作成する。「失敗するテストを書く」「テストを通す最小コードを書く」「リファクタリングする」の三つの箱を円環でつなぐ。中央に「TDD/テストファースト」と書く。各箱に「赤」「緑」「整理」という小さなラベルを日本語で添えるが、色は使わず線の種類で表現する。講義ノート風。
\end{todobox}

\subsection{構築のためのフィードバックループ}

構築で重要なのは、早く継続的なフィードバックを得ることである。アジャイル開発では、短い反復で動くソフトウェアを作り、顧客やチームから反応を得る。DevOps では、運用環境からの情報を開発へ戻すことで、実際の性能、信頼性、利用状況を学ぶ。

フィードバックループを支える技法には、自動ビルド、自動テスト、継続的統合、カナリアリリース、A/B テスト、監視、ログ分析がある。カナリアリリースは、全利用者に公開する前に、一部の利用者や環境へ変更を出して問題を観察する方法である。A/B テストは、複数の方式を比較し、実利用のデータから判断する方法である。

\section{ソフトウェア構築ツール}

構築ツールは、開発者がコードを書き、ビルドし、テストし、デバッグし、性能を測り、問題を発見するための支援を行う。ツールは生産性だけでなく品質にも影響する。ただし、ツールは判断を置き換えるものではない。ツールの結果をどう解釈し、どう改善につなげるかが重要である。

\subsection{開発環境}

開発環境、特に統合開発環境は、コード編集、コンパイル、エラー検出、ソースコード管理、ビルド、テスト、デバッグ、リファクタリング支援などを統合する。開発環境の選択は、構築効率と品質に影響する。

近年は、クラウド上の開発環境も使われる。クラウド開発環境では、ブラウザから開発でき、コンテナ化されたビルドやデプロイを利用できることがある。また、大規模言語モデルを使った AI 支援プログラミングも広がっている。AI はコード候補を生成できるが、開発者は生成されたコードをレビューし、テストし、プロジェクトへ適切に統合する責任を持つ。

\subsection{視覚的プログラミングとローコード/ゼロコードプラットフォーム}

視覚的プログラミングは、図形や部品を画面上で操作してプログラムを作る方法である。GUI ビルダーは、画面部品をドラッグアンドドロップで配置し、イベント処理のコード生成を支援する。

ローコード/ゼロコードプラットフォームは、視覚的な操作を中心にアプリケーションを作る環境である。ローコードは少量の手書きコードを必要とする。ゼロコードは、原則として手書きコードをほとんど必要としない。これらは、モデル駆動設計、視覚的プログラミング、コード生成の考え方に基づく。

このような環境は、開発速度を高める可能性がある。一方で、複雑な要件、性能、セキュリティ、拡張性、ベンダ依存、テスト容易性には注意が必要である。

\subsection{単体テストツール}

単体テストツールは、関数、クラス、モジュールなどを独立して検証するための環境を提供する。テストケースをコードとして記述し、テストランナーが自動で実行し、失敗したテストを報告する。

単体テストツールは、継続的統合と組み合わせることで価値が高まる。開発者が変更を統合するたびにテストを実行すれば、問題を早く発見できる。テストは、回帰を防ぐ安全網にもなる。

\subsection{プロファイリング、性能分析、スライシングツール}

プロファイリングツールは、実行中のプログラムを監視し、各文や関数がどれだけ実行されたか、どれだけ時間を使ったかを記録する。これにより、性能上のホットスポットを見つけられる。

プログラムスライシングは、ある変数や処理結果に影響し得る文の集合を計算する技法である。デバッグ、プログラム理解、最適化分析に役立つ。静的解析や動的解析を使ってスライスを求めるツールがある。

\begin{todobox}{構築ツールの分類図を入れる}
\textbf{画像生成用プロンプト}

白背景、モノクロ、A4 本文幅の三段構成図を作成する。上段に「開発環境」、中段に「テスト支援」、下段に「分析支援」を配置する。開発環境には「IDE」「クラウド開発環境」「AI 支援」を入れる。テスト支援には「単体テスト」「テストランナー」「継続的統合」を入れる。分析支援には「プロファイラ」「性能分析」「プログラムスライシング」を入れる。日本語ラベル、講義ノート風。
\end{todobox}

\section{トピックと参考文献の対応}

この表は、SWEBOK が示すトピック対参考資料の行列を、学習用に簡略化したものである。詳しく学ぶときは、各トピックに対応する文献を読む。

\begin{longtable}{L{0.46\linewidth}L{0.46\linewidth}}
\toprule
\textbf{トピック} & \textbf{主な参考資料} \\
\midrule
1. ソフトウェア構築の基礎 & McConnell, \textit{Code Complete}; Sommerville, \textit{Software Engineering}; Kim et al., \textit{The DevOps Handbook} \\
1.1 複雑さの最小化 & McConnell \\
1.2 変化を予期し受け入れる & McConnell; Sommerville; Kim et al. \\
1.3 検証しやすく構築する & McConnell \\
1.4 資産の再利用 & Sommerville \\
1.5 構築における標準の適用 & McConnell \\
2. 構築の管理 & McConnell; Sommerville; Kim et al. \\
2.1 ライフサイクルモデルにおける構築 & McConnell; Sommerville; Kim et al. \\
2.2 構築計画 & McConnell \\
2.3 構築測定 & McConnell \\
2.4 依存関係の管理 & Sommerville \\
3. 実務上の考慮 & McConnell; Sommerville; Heitkötter et al. \\
3.1 構築設計 & McConnell; Sommerville \\
3.2 構築言語 & McConnell \\
3.3 コーディング & McConnell \\
3.4 構築テスト & McConnell; Sommerville \\
3.5 構築における再利用 & Sommerville \\
3.6 構築品質 & McConnell; Sommerville \\
3.7 統合 & McConnell; Sommerville; Kim et al. \\
3.8 クロスプラットフォーム開発と移行 & Heitkötter et al. \\
4. 構築技術 & McConnell; Sommerville; Clements et al.; Gamma et al.; Mellor and Balcer; Null and Lobur; Silberschatz et al. \\
4.1 API の設計と利用 & Clements et al. \\
4.2 オブジェクト指向の実行時問題 & McConnell \\
4.3 パラメータ化、テンプレート、総称 & Gamma et al. \\
4.4 表明、契約による設計、防御的プログラミング & McConnell \\
4.5 エラー処理、例外処理、フォールトトレランス & McConnell \\
4.6 実行可能モデル & Mellor and Balcer \\
4.7 状態に基づく構築技法と表駆動構築技法 & McConnell \\
4.8 実行時設定と国際化 & McConnell \\
4.9 文法に基づく入力処理 & McConnell; Null and Lobur \\
4.10 並行処理の基本部品 & Silberschatz et al. \\
4.11 ミドルウェア & Clements et al.; Null and Lobur \\
4.12 分散およびクラウド基盤ソフトウェアの構築方法 & Sommerville; Kim et al. \\
4.13 異種システムの構築 & Null and Lobur \\
4.14 性能分析とチューニング & McConnell \\
4.15 プラットフォーム標準 & Heitkötter et al.; Null and Lobur; Silberschatz et al. \\
4.16 テストファーストプログラミング & McConnell; Sommerville \\
4.17 構築のためのフィードバックループ & Kim et al. \\
5. ソフトウェア構築ツール & McConnell; Sommerville \\
5.1 開発環境 & McConnell \\
5.2 視覚的プログラミングとローコード/ゼロコード & McConnell \\
5.3 単体テストツール & McConnell; Sommerville \\
5.4 プロファイリング、性能分析、スライシングツール & McConnell \\
\bottomrule
\end{longtable}

\section{発展的な読み物}

\begin{longtable}{L{0.34\linewidth}L{0.58\linewidth}}
\toprule
\textbf{文献} & \textbf{読むと得られること} \\
\midrule
McConnell, \textit{Code Complete} & コーディング、複雑さの管理、構築品質、レビュー、テスト、チューニングを実務的に学べる。構築を初めて体系的に学ぶときの中心文献である。 \\
Sommerville, \textit{Software Engineering} & ソフトウェア工学全体の中で、構築、再利用、分散システム、テストを位置づけて理解できる。 \\
Kim et al., \textit{The DevOps Handbook} & 継続的デリバリ、DevOps、フィードバックループ、運用からの学習を学べる。 \\
Fowler and Beck, \textit{Refactoring} & 既存コードの設計を改善し、変更しやすくする具体的な方法を学べる。 \\
Robert C. Martin, \textit{Clean Code} & 読みやすく保守しやすいコードを書くための原則を学べる。 \\
IEEE Std. 1517-2010 & 再利用に基づく開発、再利用のための構築、再利用による構築のプロセスを学べる。 \\
ISO/IEC 12207 & ソフトウェアライフサイクルプロセス全体の中で、構築、統合、再利用を位置づけられる。 \\
\bottomrule
\end{longtable}

\section{参考文献}

\begin{enumerate}
  \item S. McConnell, \textit{Code Complete}, 2nd edition, Microsoft Press, 2004.
  \item I. Sommerville, \textit{Software Engineering}, 10th edition, Addison-Wesley, 2016.
  \item G. Kim et al., \textit{The DevOps Handbook: How to Create World-Class Agility, Reliability and Security in Technology Organizations}, 2nd edition, IT Revolution, 2021.
  \item H. Heitkötter, S. Hanschke, and T. A. Majchrzak, ``Evaluating Cross-Platform Development Approaches for Mobile Applications,'' 2013.
  \item P. Clements et al., \textit{Documenting Software Architectures: Views and Beyond}, 2nd edition, Pearson Education, 2010.
  \item E. Gamma et al., \textit{Design Patterns: Elements of Reusable Object-Oriented Software}, Addison-Wesley Professional, 1994.
  \item S. J. Mellor and M. J. Balcer, \textit{Executable UML: A Foundation for Model-Driven Architecture}, Addison-Wesley, 2002.
  \item L. Null and J. Lobur, \textit{The Essentials of Computer Organization and Architecture}, 5th edition, Jones and Bartlett Publishers, 2018.
  \item A. Silberschatz et al., \textit{Operating System Concepts}, 8th edition, Wiley, 2008.
  \item IEEE Std. 1517-2010, \textit{IEEE Standard for Information Technology - Software Life Cycle Processes - Reuse Processes}.
  \item ISO/IEC 12207, \textit{Systems and Software Engineering - Software Life Cycle Processes}.
  \item M. Fowler and K. Beck, \textit{Refactoring: Improving the Design of Existing Code}, 2nd edition, Addison-Wesley.
  \item R. C. Martin, \textit{Clean Code: A Handbook of Agile Software Craftsmanship}, Pearson Education.
\end{enumerate}

\section{画像 TODO 一覧}

本文に配置した画像 TODO を再掲する。実際に図を作る場合は、各 TODO のプロンプトをそのまま画像生成に使うと、A4 資料に合う図を作りやすい。

\begin{enumerate}
  \item 第04章の全体地図。
  \item 構築と他知識領域の関係図。
  \item 構築の五つの基礎概念図。
  \item ライフサイクル別の構築位置比較図。
  \item 依存関係の供給網図。
  \item 構築言語の分類図。
  \item 段階的統合と継続的統合の流れ。
  \item API ファーストの流れ。
  \item 表明・契約・防御的プログラミングの関係図。
  \item 実行可能モデル、PIM、PSM の関係図。
  \item 状態機械と表駆動の比較図。
  \item 分散・クラウド構築の概念図。
  \item 性能分析とチューニングの流れ。
  \item テスト駆動開発のサイクル図。
  \item 構築ツールの分類図。
\end{enumerate}

\section{章末確認問題}

\begin{enumerate}
  \item ソフトウェア構築は、単なるコーディングとどう違うか。
  \item 複雑さを最小化することが、なぜテストや保守に影響するのか。
  \item 「変化を予期する」と「過剰に柔軟に作る」は何が違うか。
  \item 検証しやすく構築するために、コード構造で注意すべきことを三つ挙げよ。
  \item 再利用のための構築と、再利用による構築の違いを説明せよ。
  \item 構築計画で、部品の統合順序を決めることが重要な理由を説明せよ。
  \item 依存関係管理で、ライセンスと脆弱性を確認すべき理由を説明せよ。
  \item 単体テストと統合テストの違いを説明せよ。
  \item API 設計で「誤用しにくい」ことが重要な理由を説明せよ。
  \item 表明、契約による設計、防御的プログラミングの違いを説明せよ。
  \item 分散・クラウド基盤ソフトウェアで、単一マシンのソフトウェアより難しくなる点を三つ挙げよ。
  \item 性能チューニングで、測定なしに最適化してはならない理由を説明せよ。
\end{enumerate}

\begin{plainbox}
\textbf{解答の方向性}

1 は、構築がコーディング、検証、単体テスト、統合テスト、デバッグ、再利用、統合、依存関係管理などを含む点を述べる。2 は、複雑なコードほど実行経路が多く、理解とテストが難しくなる点を述べる。3 は、予期された変化への備えと、根拠のない過剰設計の違いを述べる。4 以降は、本文の「単純で読みやすいコード」「早いフィードバック」「再利用のリスク」「段階的統合」「契約」「測定に基づく改善」といった語を使って説明できる。
\end{plainbox}

\end{document}
